% GENERATED FILE. Edit canonical lessons, then run npm run generate:books.
% canonical-source-hash: fnv1a64:dc73e72f6cd48b3b
% canonical-lessons: AR-C411-ajaza, AR-C411-khadaa, AR-C411-zalama, AR-C411-kafaha, AR-C411-aqarra

\chapter{Rules and wrongs}
\label{ch:ar-rules-and-wrongs}
\hlchaptermodality{411}

\begin{chapteropening}
\textbf{By the end of this chapter:} \emph{I can say that someone authorized, was subject to, wronged, struggled against or formally approved something.}

Complete the last lesson of chapter 411: I can say that someone authorized, was subject to, wronged, struggled against or formally approved something.
\end{chapteropening}

\section[\texorpdfstring{\ar{أجاز} (ajāza)}{أجاز (ajāza)}]{\ar{أجاز} (ajāza) — to authorize}

\label{lesson:AR-C411-ajaza}



\begin{quote}
Before the new one: say the Arabic for to comply, then the Arabic for to exempt.
\end{quote}

\subsection*{You'll want to know: \ar{أجاز}}
\textbf{\ar{أجاز}} — \emph{ajāza} — \textquotedblleft{}to authorize\textquotedblright{}.

\begin{quote}
\ar{أجاز} \ar{الأستاذ} \ar{استعمال} \ar{القاموس}.
\end{quote}

\begin{quote}
\emph{ajāza al-ustādh istiʿmāl al-qāmūs.}
\end{quote}

\begin{quote}
The professor authorized the use of the dictionary.
\end{quote}

Present: \ar{يجيز}, \emph{yujīz}.

\begin{cousinweb}
\textbf{\ar{أجاز}} (\emph{ajāza}) is built on the native Arabic root \emph{j-w-z}, \textquotedblleft{}to pass, to go beyond\textquotedblright{} (Lane). It is Form IV, to let a thing pass: to allow, to authorize.
\end{cousinweb}

\begin{grammarlens}[title={What you've built}]
One more word for saying what happens if.
\end{grammarlens}

\subsection*{Your turn}
\begin{itemize}
\raggedright
  \item \emph{Say it:} \emph{ajāza}
  \item \emph{Say it:} \emph{ajāza}, once more
  \item \emph{Say it:} \emph{ajāza al-ustādh istiʿmāl al-qāmūs.} — The professor authorized the use of the dictionary.
  \item \emph{Recall:} say the Arabic for to comply, then the Arabic for to exempt, then say \emph{ajāza} again
\end{itemize}

\begin{tcolorbox}[breakable,colback=teal!4,colframe=teal!35!black,title={Before you move on}]
What does \ar{أجاز} mean? (\textquotedblleft{}To authorize\textquotedblright{}.) Say it once more.
\end{tcolorbox}
\section[\texorpdfstring{\ar{خضع} (khaḍaʿa)}{خضع (khaḍaʿa)}]{\ar{خضع} (khaḍaʿa) — to be subject to}

\label{lesson:AR-C411-khadaa}



\begin{quote}
Before the new one: say the Arabic for to exempt, then the Arabic for to authorize.
\end{quote}

\subsection*{You'll want to know: \ar{خضع}}
\textbf{\ar{خضع}} — \emph{khaḍaʿa} — \textquotedblleft{}to be subject to\textquotedblright{}.

\begin{quote}
\ar{يخضع} \ar{الجميع} \ar{للقانون}.
\end{quote}

\begin{quote}
\emph{yakhḍaʿu al-jamīʿ li-l-qānūn.}
\end{quote}

\begin{quote}
Everyone is subject to the law.
\end{quote}

Present: \ar{يخضع}, \emph{yakhḍaʿ}.

It takes \ar{لـ}: \ar{خضع} \ar{لـ} …. \ar{خضع} \ar{لعملية}, underwent an operation.

\begin{cousinweb}
\textbf{\ar{خضع}} (\emph{khaḍaʿa}) is built on the native Arabic root \emph{kh-ḍ-ʿ}, \textquotedblleft{}to be lowly, submissive\textquotedblright{} (Lane). \emph{Khaḍaʿa} is a Form I verb: to submit, to be subject to.
\end{cousinweb}

\begin{grammarlens}[title={What you've built}]
One more word for saying what happens if.
\end{grammarlens}

\subsection*{Your turn}
\begin{itemize}
\raggedright
  \item \emph{Say it:} \emph{khaḍaʿa}
  \item \emph{Say it:} \emph{khaḍaʿa}, once more
  \item \emph{Say it:} \emph{yakhḍaʿu al-jamīʿ li-l-qānūn.} — Everyone is subject to the law.
  \item \emph{Recall:} say the Arabic for to exempt, then the Arabic for to authorize, then say \emph{khaḍaʿa} again
\end{itemize}

\begin{tcolorbox}[breakable,colback=teal!4,colframe=teal!35!black,title={Before you move on}]
What does \ar{خضع} mean? (\textquotedblleft{}To be subject to\textquotedblright{}.) Say it once more.
\end{tcolorbox}
\section[\texorpdfstring{\ar{ظلم} (ẓalama)}{ظلم (ẓalama)}]{\ar{ظلم} (ẓalama) — to wrong}

\label{lesson:AR-C411-zalama}



\begin{quote}
Before the new one: say the Arabic for to authorize, then the Arabic for to be subject to.
\end{quote}

\subsection*{You'll want to know: \ar{ظلم}}
\textbf{\ar{ظلم}} — \emph{ẓalama} — \textquotedblleft{}to wrong\textquotedblright{}.

\begin{quote}
\ar{ظلم} \ar{المدير} \ar{العامل} \ar{الجديد}.
\end{quote}

\begin{quote}
\emph{ẓalama al-mudīr al-ʿāmil al-jadīd.}
\end{quote}

\begin{quote}
The manager wronged the new worker.
\end{quote}

Present: \ar{يظلم}, \emph{yaẓlim}.

\begin{cousinweb}
\textbf{\ar{ظلم}} (\emph{ẓalama}) is built on the native Arabic root \emph{ẓ-l-m}, \textquotedblleft{}to put a thing in the wrong place\textquotedblright{} (Lane). \emph{Ẓalama} is a Form I verb: to do someone wrong. \ar{ظالم}, unjust, is its active participle.
\end{cousinweb}

\begin{grammarlens}[title={What you've built}]
One more word for saying what happens if.
\end{grammarlens}

\subsection*{Your turn}
\begin{itemize}
\raggedright
  \item \emph{Say it:} \emph{ẓalama}
  \item \emph{Say it:} \emph{ẓalama}, once more
  \item \emph{Say it:} \emph{ẓalama al-mudīr al-ʿāmil al-jadīd.} — The manager wronged the new worker.
  \item \emph{Recall:} say the Arabic for to authorize, then the Arabic for to be subject to, then say \emph{ẓalama} again
\end{itemize}

\begin{tcolorbox}[breakable,colback=teal!4,colframe=teal!35!black,title={Before you move on}]
What does \ar{ظلم} mean? (\textquotedblleft{}To wrong\textquotedblright{}.) Say it once more.
\end{tcolorbox}
\section[\texorpdfstring{\ar{كافح} (kāfaḥa)}{كافح (kāfaḥa)}]{\ar{كافح} (kāfaḥa) — to struggle against}

\label{lesson:AR-C411-kafaha}



\begin{quote}
Before the new one: say the Arabic for to be subject to, then the Arabic for to wrong.
\end{quote}

\subsection*{You'll want to know: \ar{كافح}}
\textbf{\ar{كافح}} — \emph{kāfaḥa} — \textquotedblleft{}to struggle against\textquotedblright{}.

\begin{quote}
\ar{تكافح} \ar{الدولة} \ar{الفقر}.
\end{quote}

\begin{quote}
\emph{tukāfiḥu ad-dawla al-faqr.}
\end{quote}

\begin{quote}
The state is fighting poverty.
\end{quote}

Present: \ar{يكافح}, \emph{yukāfiḥ}.

What is fought is the direct object: \ar{كافح} \ar{الفساد}, fought corruption.

\begin{cousinweb}
\textbf{\ar{كافح}} (\emph{kāfaḥa}) is built on the native Arabic root \emph{k-f-ḥ}, \textquotedblleft{}to meet face to face\textquotedblright{} (Lane). It is Form III, to face an enemy head on: to struggle against.
\end{cousinweb}

\begin{grammarlens}[title={What you've built}]
One more word for saying what happens if.
\end{grammarlens}

\subsection*{Your turn}
\begin{itemize}
\raggedright
  \item \emph{Say it:} \emph{kāfaḥa}
  \item \emph{Say it:} \emph{kāfaḥa}, once more
  \item \emph{Say it:} \emph{tukāfiḥu ad-dawla al-faqr.} — The state is fighting poverty.
  \item \emph{Recall:} say the Arabic for to be subject to, then the Arabic for to wrong, then say \emph{kāfaḥa} again
\end{itemize}

\begin{tcolorbox}[breakable,colback=teal!4,colframe=teal!35!black,title={Before you move on}]
What does \ar{كافح} mean? (\textquotedblleft{}To struggle against\textquotedblright{}.) Say it once more.
\end{tcolorbox}
\section[\texorpdfstring{\ar{أقر} (aqarra)}{أقر (aqarra)}]{\ar{أقر} (aqarra) — to approve formally}

\label{lesson:AR-C411-aqarra}



\begin{quote}
Before the new one: say the Arabic for to wrong, then the Arabic for to struggle against.
\end{quote}

\subsection*{You'll want to know: \ar{أقر}}
\textbf{\ar{أقر}} — \emph{aqarra} — \textquotedblleft{}to approve formally\textquotedblright{}.

\begin{quote}
\ar{أقر} \ar{المجلس} \ar{الخطة} \ar{الجديدة}.
\end{quote}

\begin{quote}
\emph{aqarra al-majlis al-khiṭṭa al-jadīda.}
\end{quote}

\begin{quote}
The council formally approved the new plan.
\end{quote}

Present: \ar{يقر}, \emph{yuqirr}.

Also to admit: \ar{أقر} \ar{بالخطأ}, admitted the mistake.

\begin{cousinweb}
\textbf{\ar{أقر}} (\emph{aqarra}) is built on the native Arabic root \emph{q-r-r}, \textquotedblleft{}to settle, to be still\textquotedblright{} (Lane). It is Form IV, to make a thing settled: to approve it formally, or to admit it.
\end{cousinweb}

\begin{grammarlens}[title={What you've built}]
One more word for saying what happens if.
\end{grammarlens}

\subsection*{Your turn}
\begin{itemize}
\raggedright
  \item \emph{Say it:} \emph{aqarra}
  \item \emph{Say it:} \emph{aqarra}, once more
  \item \emph{Say it:} \emph{aqarra al-majlis al-khiṭṭa al-jadīda.} — The council formally approved the new plan.
  \item \emph{Recall:} say the Arabic for to wrong, then the Arabic for to struggle against, then say \emph{aqarra} again
\end{itemize}

\begin{tcolorbox}[breakable,colback=teal!4,colframe=teal!35!black,title={Before you move on}]
What does \ar{أقر} mean? (\textquotedblleft{}To approve formally\textquotedblright{}.) Say it once more.
\end{tcolorbox}
